\clearpage\thispagestyle{empty}
\pdfbookmark[0]{Quellen zu Komet (10. Februar 1821)}{cometfeb1821-sources}
\begin{center}\textbf{Quellen zu Komet (10. Februar 1821)}\end{center}
\noindent Diese KI-Übertragung ins heutige Deutsch folgt \textit{Comet} in den \href{https://gdz.sub.uni-goettingen.de/id/PPN236007467}{\textit{Werken}, Band~VI} (Göttingen, 1874), S.~435. Der \href{https://images.p.sub.uni-goettingen.de/iiif/image/gdz:PPN319721507_1821:00000217/full/1600,/0/default.jpg}{Erstdruck} erschien unter \textit{Göttingen} in den \textit{Göttingischen gelehrten Anzeigen}, Nr.~23, 10. Februar 1821, S.~217, im Exemplar der Niedersächsischen Staats- und Universitätsbibliothek Göttingen. Er nennt Nicolet als Entdecker und Gauß für das Auffinden und die Beobachtung des Kometen in Göttingen.
\par\smallskip\noindent Der Erstdruck bezeichnet die Minuten und Sekunden der Beobachtungszeit mit Strichen. Die Werkausgabe verwendet hochgestellte Buchstaben. In der Deklinationszeile fehlen in beiden Ausgaben die Einheitenzeichen. KI-Übertragung und KI-Satz.
\clearpage
